% GENERATED FILE. Edit canonical lessons, then run npm run generate:books.
% canonical-source-hash: fnv1a64:51f70ac18c22201e
% canonical-lessons: BN-C448-biroho, BN-C448-akangkha, BN-C448-udas, BN-C448-bekul, BN-C448-hahakar-kora

\chapter{Longing}
\label{ch:bn-longing}
\hlchaptermodality{448}

\begin{chapteropening}
\textbf{By the end of this chapter:} \emph{I can talk about missing someone and longing for something.}

Complete the last lesson of chapter 448: I can talk about missing someone and longing for something.
\end{chapteropening}

\section[\texorpdfstring{\bn{বিরহ} (birôhô)}{বিরহ (birôhô)}]{\bn{বিরহ} (birôhô) — the pain of separation; longing for someone far away}

\label{lesson:BN-C448-biroho}



\begin{quote}
Before the new one: say the Bengali for to rue, then the Bengali for to admit.
\end{quote}

\subsection*{You'll want to know: \bn{বিরহ}}
\textbf{\bn{বিরহ}} — \emph{birôhô} — \textquotedblleft{}the pain of separation; longing for someone far away\textquotedblright{}.

\begin{quote}
\bn{এই} \bn{গানটা} \bn{বিরহের} \bn{গান।}
\end{quote}

\begin{quote}
\emph{ei gānṭā birôher gān.}
\end{quote}

\begin{quote}
This is a song of longing for someone far away.
\end{quote}

\begin{cousinweb}
\textbf{\bn{বিরহ}} (\emph{birôhô}) is a tatsama, Sanskrit \emph{viraha}, \textquotedblleft{}parting, separation\textquotedblright{}, especially of lovers (Monier-Williams); a great deal of Bengali song is built on it.
\end{cousinweb}

\begin{grammarlens}[title={What you've built}]
One more word for describing what you lived and felt.
\end{grammarlens}

\subsection*{Your turn}
\begin{itemize}
\raggedright
  \item \emph{Say it:} \emph{birôhô}
  \item \emph{Say it:} \emph{birôhô}, once more
  \item \emph{Say it:} \emph{ei gānṭā birôher gān.} — This is a song of longing for someone far away.
  \item \emph{Recall:} say the Bengali for to rue, then the Bengali for to admit, then say \emph{birôhô} again
\end{itemize}

\begin{tcolorbox}[breakable,colback=teal!4,colframe=teal!35!black,title={Before you move on}]
What does \bn{বিরহ} mean? (\textquotedblleft{}The pain of separation; longing for someone far away\textquotedblright{}.) Say it once more.
\end{tcolorbox}
\section[\texorpdfstring{\bn{আকাঙ্ক্ষা} (ākāngkhā)}{আকাঙ্ক্ষা (ākāngkhā)}]{\bn{আকাঙ্ক্ষা} (ākāngkhā) — a yearning; a deep desire}

\label{lesson:BN-C448-akangkha}



\begin{quote}
Before the new one: say the Bengali for to admit, then the Bengali for the pain of separation.
\end{quote}

\subsection*{You'll want to know: \bn{আকাঙ্ক্ষা}}
\textbf{\bn{আকাঙ্ক্ষা}} — \emph{ākāngkhā} — \textquotedblleft{}a yearning; a deep desire\textquotedblright{}.

\begin{quote}
\bn{বড়} \bn{শিল্পী} \bn{হওয়াই} \bn{মীরার} \bn{আকাঙ্ক্ষা।}
\end{quote}

\begin{quote}
\emph{bôṛo shilpī hôoyāi mīrār ākāngkhā.}
\end{quote}

\begin{quote}
To become a great artist is Mira's deepest desire.
\end{quote}

\begin{cousinweb}
\textbf{\bn{আকাঙ্ক্ষা}} (\emph{ākāngkhā}) is a tatsama, Sanskrit \emph{ākāṅkṣā}, \textquotedblleft{}desire, wish\textquotedblright{} (Monier-Williams).
\end{cousinweb}

\begin{grammarlens}[title={What you've built}]
One more word for describing what you lived and felt.
\end{grammarlens}

\subsection*{Your turn}
\begin{itemize}
\raggedright
  \item \emph{Say it:} \emph{ākāngkhā}
  \item \emph{Say it:} \emph{ākāngkhā}, once more
  \item \emph{Say it:} \emph{bôṛo shilpī hôoyāi mīrār ākāngkhā.} — To become a great artist is Mira's deepest desire.
  \item \emph{Recall:} say the Bengali for to admit, then the Bengali for the pain of separation, then say \emph{ākāngkhā} again
\end{itemize}

\begin{tcolorbox}[breakable,colback=teal!4,colframe=teal!35!black,title={Before you move on}]
What does \bn{আকাঙ্ক্ষা} mean? (\textquotedblleft{}A yearning; a deep desire\textquotedblright{}.) Say it once more.
\end{tcolorbox}
\section[\texorpdfstring{\bn{উদাস} (udās)}{উদাস (udās)}]{\bn{উদাস} (udās) — wistful; far away in one's thoughts}

\label{lesson:BN-C448-udas}



\begin{quote}
Before the new one: say the Bengali for the pain of separation, then the Bengali for a yearning.
\end{quote}

\subsection*{You'll want to know: \bn{উদাস}}
\textbf{\bn{উদাস}} — \emph{udās} — \textquotedblleft{}wistful; far away in one's thoughts\textquotedblright{}.

\begin{quote}
\bn{জানালার} \bn{পাশে} \bn{বসে} \bn{মীরা} \bn{উদাস} \bn{হয়ে} \bn{আকাশ} \bn{দেখছিল।}
\end{quote}

\begin{quote}
\emph{jānālār pāshe bôse mīrā udās hôye ākāsh dekhchhilô.}
\end{quote}

\begin{quote}
Sitting by the window, Mira gazed wistfully at the sky.
\end{quote}

\begin{cousinweb}
\textbf{\bn{উদাস}} (\emph{udās}) is a tatsama, Sanskrit \emph{udāsa}, \textquotedblleft{}indifference, apathy\textquotedblright{} (Monier-Williams); in Bengali it is the far-away look of someone lost in thought.
\end{cousinweb}

\begin{grammarlens}[title={What you've built}]
One more word for describing what you lived and felt.
\end{grammarlens}

\subsection*{Your turn}
\begin{itemize}
\raggedright
  \item \emph{Say it:} \emph{udās}
  \item \emph{Say it:} \emph{udās}, once more
  \item \emph{Say it:} \emph{jānālār pāshe bôse mīrā udās hôye ākāsh dekhchhilô.} — Sitting by the window, Mira gazed wistfully at the sky.
  \item \emph{Recall:} say the Bengali for the pain of separation, then the Bengali for a yearning, then say \emph{udās} again
\end{itemize}

\begin{tcolorbox}[breakable,colback=teal!4,colframe=teal!35!black,title={Before you move on}]
What does \bn{উদাস} mean? (\textquotedblleft{}Wistful; far away in one's thoughts\textquotedblright{}.) Say it once more.
\end{tcolorbox}
\section[\texorpdfstring{\bn{ব্যাকুল} (bækul)}{ব্যাকুল (bækul)}]{\bn{ব্যাকুল} (bækul) — anxious with longing; eager}

\label{lesson:BN-C448-bekul}



\begin{quote}
Before the new one: say the Bengali for a yearning, then the Bengali for wistful.
\end{quote}

\subsection*{You'll want to know: \bn{ব্যাকুল}}
\textbf{\bn{ব্যাকুল}} — \emph{bækul} — \textquotedblleft{}anxious with longing; eager\textquotedblright{}.

\begin{quote}
\bn{ছেলের} \bn{ফোনের} \bn{জন্য} \bn{বাবা} \bn{ব্যাকুল} \bn{হয়ে} \bn{আছেন।}
\end{quote}

\begin{quote}
\emph{chheler phoner jônyô bābā bækul hôye āchhen.}
\end{quote}

\begin{quote}
Father is waiting anxiously for a call from his son.
\end{quote}

\begin{cousinweb}
\textbf{\bn{ব্যাকুল}} (\emph{bækul}) is a tatsama, Sanskrit \emph{vyākula}, \textquotedblleft{}filled with\textquotedblright{} and \textquotedblleft{}troubled, perplexed\textquotedblright{} (Monier-Williams).
\end{cousinweb}

\begin{grammarlens}[title={What you've built}]
One more word for describing what you lived and felt.
\end{grammarlens}

\subsection*{Your turn}
\begin{itemize}
\raggedright
  \item \emph{Say it:} \emph{bækul}
  \item \emph{Say it:} \emph{bækul}, once more
  \item \emph{Say it:} \emph{chheler phoner jônyô bābā bækul hôye āchhen.} — Father is waiting anxiously for a call from his son.
  \item \emph{Recall:} say the Bengali for a yearning, then the Bengali for wistful, then say \emph{bækul} again
\end{itemize}

\begin{tcolorbox}[breakable,colback=teal!4,colframe=teal!35!black,title={Before you move on}]
What does \bn{ব্যাকুল} mean? (\textquotedblleft{}Anxious with longing; eager\textquotedblright{}.) Say it once more.
\end{tcolorbox}
\section[\texorpdfstring{\bn{হাহাকার} \bn{করা} (hāhākār kôrā)}{হাহাকার করা (hāhākār kôrā)}]{\bn{হাহাকার} \bn{করা} (hāhākār kôrā) — to cry out in despair}

\label{lesson:BN-C448-hahakar-kora}



\begin{quote}
Before the new one: say the Bengali for wistful, then the Bengali for anxious with longing.
\end{quote}

\subsection*{You'll want to know: \bn{হাহাকার} \bn{করা}}
\textbf{\bn{হাহাকার} \bn{করা}} — \emph{hāhākār kôrā} — \textquotedblleft{}to cry out in despair\textquotedblright{}.

\begin{quote}
\bn{বন্যার} \bn{পরে} \bn{গ্রামের} \bn{লোকেরা} \bn{হাহাকার} \bn{করছিল।}
\end{quote}

\begin{quote}
\emph{bônyār pôre grāmer lokerā hāhākār kôrchhilô.}
\end{quote}

\begin{quote}
After the flood the villagers were crying out in despair.
\end{quote}

Past (he/she …): \bn{তিনি} \bn{ছেলের} \bn{জন্য} \bn{হাহাকার} \bn{করলেন}, \emph{tini chheler jônyô hāhākār kôrlen}.

\begin{cousinweb}
\textbf{\bn{হাহাকার} \bn{করা}} (\emph{hāhākār kôrā}) is \emph{hāhākār}, the Sanskrit cry \emph{hā}, \textquotedblleft{}alas\textquotedblright{}, said twice, with \emph{kāra}, \textquotedblleft{}making\textquotedblright{} (Monier-Williams), and \emph{kôrā}, \textquotedblleft{}to do\textquotedblright{} (CDIAL 2814).
\end{cousinweb}

\begin{grammarlens}[title={What you've built}]
One more word for describing what you lived and felt.
\end{grammarlens}

\subsection*{Your turn}
\begin{itemize}
\raggedright
  \item \emph{Say it:} \emph{hāhākār kôrā}
  \item \emph{Say it:} \emph{hāhākār kôrā}, once more
  \item \emph{Say it:} \emph{bônyār pôre grāmer lokerā hāhākār kôrchhilô.} — After the flood the villagers were crying out in despair.
  \item \emph{Recall:} say the Bengali for wistful, then the Bengali for anxious with longing, then say \emph{hāhākār kôrā} again
\end{itemize}

\begin{tcolorbox}[breakable,colback=teal!4,colframe=teal!35!black,title={Before you move on}]
What does \bn{হাহাকার} \bn{করা} mean? (\textquotedblleft{}To cry out in despair\textquotedblright{}.) Say it once more.
\end{tcolorbox}
